\section{Introduction}

Fix an integer \(d\ge2\), a compact convex set \(K\subset\R^d\) with
nonempty interior, and the uniform probability measure
\[
 \rho=\frac{1_K}{|K|}\,dx.
\]
For a nonempty compact set \(Y\subset\R^d\), let \(\Prob(Y)\) denote the
Borel probability measures supported on \(Y\). The quadratic optimal
transport map from \(\rho\) to \(\mu\in\Prob(Y)\) is denoted by \(T_\mu\).
It minimizes
\[
 \int_K |x-T(x)|^2\,d\rho(x)
 \quad\text{subject to }T_\#\rho=\mu
\]
and is unique up to a \(\rho\)-null set. Such maps are gradients of convex
potentials, as in Brenier's theorem \cite{Brenier1991}. We give direct
existence and uniqueness arguments for the maps used here.

For probability measures \(\mu,\nu\) with compact support, write
\[
 W_2(\mu,\nu)
 =\inf_{\pi\in\Pi(\mu,\nu)}
   \left(\int_{\R^d\times\R^d}|y-z|^2\,d\pi(y,z)\right)^{1/2},
\]
where \(\Pi(\mu,\nu)\) is the set of probability couplings. We use
\(W_1(\mu,\nu)=\inf_{\pi\in\Pi(\mu,\nu)}\int|y-z|\,d\pi\)
only when comparing earlier bounds. Our result is the following uniform
estimate.

\begin{theorem}\label{main:stability}
For every \(K,Y\) as above there is a finite constant \(C(K,Y)\) such that
\[
 \|T_\mu-T_\nu\|_{L^2(\rho)}
 \le C(K,Y)\,W_2(\mu,\nu)^{1/3}
 \qquad\bigl(\mu,\nu\in\Prob(Y)\bigr).
\]
The targets may be atomic, singular, or absolutely continuous. The
exponent \(1/3\) cannot be replaced by any larger exponent in general,
even for \(K=Y=[-1,1]^d\) and targets with exactly three atoms.
\end{theorem}

The constant is independent of atom count, positive atom masses, site
separation, and density bounds. Neither smoothness nor strict convexity
of \(K\) is required, and \(Y\) need not be convex. An explicit admissible
constant is obtained in \Cref{gradient:finite}.

\subsection{History and significance}

Quantitative stability asks how strongly the representation of a target
measure by its optimal map depends on the target. Earlier contributions
include Ambrosio's local one-half H\"older estimate along target curves
near regular data, reported by Gigli
\cite[Corollary~3.4]{Gigli2011}, and Berman's estimates for varying
sources against a fixed target whose density is bounded below on a
convex domain \cite{Berman2021}. M\'erigot, Delalande, and Chazal
used dualization to obtain target-uniform bounds from Berman's work,
with dimension-dependent exponents, and then proved a
dimension-independent \(W_1^{2/15}\) bound for a uniform convex source
\cite[Corollary~2.4 and Theorem~3.1]{MerigotDelalandeChazal2020}.

Potential estimates followed by convex-gradient interpolation became
central to the target-uniform problem. Delalande and M\'erigot proved
an \(L^2\)-map bound of order \(W_1^{1/6}\) for source densities bounded
above and below on a convex domain and targets in a fixed compact set
\cite[Theorem~4.2]{DelalandeMerigot2023}. Their interpolation inequality
\cite[Proposition~4.1]{DelalandeMerigot2023} is the antecedent of our
last step; the improvement here is the linear \(W_2\) estimate for
potential differences. Letrouit and M\'erigot extended potential
stability to John domains by local-to-global variance arguments, and
obtained the corresponding map estimate when the boundary is
rectifiable with finite \((d-1)\)-dimensional measure
\cite[Theorem~1.7]{LetrouitMerigot2024}.

Divol, Niles-Weed, and Pooladian obtained a
\(W_2^{1/3}\)-map estimate for regular source densities and
nondegenerate finite targets
\cite[Theorem~4.3]{DivolNilesWeedPooladian2025}. Its constant depends on
the number of atoms, a positive lower bound for their masses, and
separation and angular geometry, whereas \Cref{main:stability} has one
constant for all Borel targets, including degenerating atomic measures.
At the other end of the regularity scale, Lipschitz \(L^2\) stability
in \(W_2\) holds for uniformly positive H\"older-continuous densities
on suitably regular uniformly convex domains
\cite[Theorem~1.1]{CajaLopezDelgadinoKitagawa2026}. Its constants depend
on nondegeneracy and regularity conditions that exclude our atomic
targets.

Cazelles, Pauwels, and Portales obtained a \(W_1^{1/5}\) map bound for
bounded source densities and compactly supported marginals in their
analysis of Monge-map estimation
\cite[Corollary~4.9]{CazellesPauwelsPortales2026}.
Han and Zhu \cite[Proposition~2.7]{HanZhu2026} report a sharp
\(W_1^{1/4}\) estimate of M\'erigot \cite{Merigot2026}; see also
\cite[Section~2.3]{ChewiEichingerPooladian2026}. Since \(W_1\le W_2\),
a quarter-power estimate in \(W_1\) gives at most a quarter-power
estimate in \(W_2\) by that comparison alone. In the fixed-cube family
of \Cref{sharpness:exponent}, \(W_1\) and \(W_2\) have different scaling,
so its \(W_2^{1/3}\) obstruction is compatible with a sharp
\(W_1^{1/4}\) estimate.

Letrouit proved an obstruction to exponents above \(1/3\) for a
nonconvex source and formulated uniform square-root stability in \(W_2\)
for a uniform convex source as a conjecture
\cite[Theorem~2 and Conjecture~3]{Letrouit2026}. The fixed-cube
three-plane construction proved in \Cref{sharpness:exponent} disproves
that conjectured exponent and every other exponent above \(1/3\), even
with only three target atoms. The positive bound in \Cref{main:stability}
identifies the matching uniform endpoint. Together, the two parts
determine the sharp exponent within the stated uniform-source class.

\subsection{Sharpness and the finite-cell argument}

We begin in \Cref{sec:sharpness} with the obstruction. A maximum of
three affine functions has a thin central cell between two outer cells.
Tilting the middle affine function preserves its cell mass but switches
source points between slopes at order-one separation. The resulting
three-atom targets give the sharpness assertion. The remaining sections
prove the matching upper bound.

For distinct sites \(y_1,\ldots,y_N\) and intercepts
\(c=(c_1,\ldots,c_N)\), set
\[
 u_c(x)=\max_i\{x\cdot y_i+c_i\},\qquad
 P_i(c)=\{x\in K:u_c(x)=x\cdot y_i+c_i\}.
\]
The cells partition \(K\) up to null sets, and \(\nabla u_c=y_i\)
almost everywhere on cell \(i\). Write
\(M_i(c)=\rho(P_i(c))\) and
\(B_i(c)=\int_{P_i(c)}x\,d\rho(x)\) for its mass and first moment.
\Cref{sec:cells} establishes the cell calculus and shows that the
intercepts realizing prescribed positive masses depend differentiably
on the sites. At a configuration with all \(M_i>0\), fix
\(q\in\R^N\) and let dots denote derivatives along \(c+s q\) at
\(s=0\), with the sites fixed.
The central estimate, proved in \Cref{sec:moments}, is
\[
 \sum_i\frac{|\dot B_i|^2}{M_i}
 \le C_0R^2\sum_i\frac{|\dot M_i|^2}{M_i},
 \qquad K\subset B(0,R),
\]
with an absolute constant \(C_0\). Couple the uniform laws on the cells
at \(c-s q\) and \(c+s q\). Their midpoint law, multiplied by the
geometric mean of their masses, can be dominated by the source measure
on the cell at \(c\). We prove this density bound by Knothe's triangular
rearrangement \cite{Knothe1957}; it is also a consequence of displacement
concavity \cite[Theorem~2.3]{McCann1997}. Testing with
\(R^2-|x|^2\) and summing over the partitions controls displacement by
mass changes. This produces the moment estimate with a constant that
remains uniform as masses or faces become small.

In \Cref{sec:finite-stability}, we move the sites of two finite targets
along an optimal target coupling while keeping the coupling masses
fixed. The derivative of the cell-mass constraint is a graph-Laplacian
equation, familiar from semi-discrete transport
\cite{KitagawaMerigotThibert2019}. We test it against an intercept
direction whose mass derivative is the weighted intercept velocity.
The moment inequality then bounds that velocity directly in its
weighted \(L^2\) norm. This avoids any quantitative bound on the inverse
Laplacian and yields an \(L^2(\rho)\) potential difference of order
\(W_2(\mu,\nu)\).

Finally, \Cref{sec:gradient} uses coordinate slices to bound the squared
gradient difference by a constant times
\[
 h^{-2}\|u_0-u_1\|_{L^2(\rho)}^2+h.
\]
The potential estimate and the choice \(h=W_2(\mu,\nu)^{2/3}\) give
the one-third exponent. Finite approximation and supporting inequalities
then extend the estimate to all Borel targets. The source's convexity
enters through its convex cells and interval slices; no smoothness of
its boundary is used.

\subsection{Conventions}

Lebesgue volume is denoted by \(|\cdot|\), and surface measure on an affine
hyperplane by \(dS\). All source masses and moments are normalized by
\(|K|\). We choose fixed \(R,L\ge1\) such that
\(K\subset B(0,R)\) and \(Y\subset B(0,L)\). A finite maximum of affine
functions with slopes in \(B(0,L)\) is convex and \(L\)-Lipschitz.
Ties between distinct slopes occur on a finite union of hyperplanes and
therefore have zero source measure.
